\documentclass[11pt]{ctexart}
\usepackage[a4paper,margin=1in]{geometry}
\usepackage{graphicx}
\usepackage{xcolor}
\usepackage{hyperref}
\definecolor{refgray}{gray}{0.45}
\title{\textbf{Market Views｜全球投行叙事汇编}\\[0.4em]{\large AI需求从算力向存储与光互连扩散，中国财政紧缩与贸易结构重塑主导宏观定价分歧}}
\author{KC桌面 自动生成}
\date{260727}
\begin{document}
\maketitle
\tableofcontents
\newpage
\section{一页摘要}
\begin{itemize}
  \item AI基础设施投资正从GPU算力向存储、光互连和功率半导体等配套领域扩散，代理型AI预计推动全球存储用量在2026-2030年间增长超七倍，供需缺口持续存在。
  \item 内存接口芯片市场中，MRDIMM渗透率假设分歧导致2030年TAM预期差异巨大（200亿美元 vs 80亿美元），存储供给瓶颈的持续性也存在产能扩张与多重瓶颈的争论。
  \item 中国互联网平台监管进入终局阶段，携程反垄断处罚落地消除不确定性，但排他协议取消后竞争格局重塑开启，市场对利润率长期影响存在分歧。
  \item 全球奢侈品欧洲市场表现分化由三个结构性因素决定：对中国游客依赖度、区域价差大小、美国游客消费替代能力，硬奢与软奢表现分化显著。
  \item 中国二季度GDP增速放缓至4.3\%，主因广义财政赤字率从11.0\%大幅收窄至7.0\%，土地出让收入同比骤降42\%是核心拖累。
  \item 中国贸易模式转向“多出少进”，对非美发达市场出口增长约20\%，已使当地商品价格平均下降0.7\%，日本和欧元区受影响最大。
  \item 美联储7月会议路径分歧加剧：通胀改善支持按兵不动，但地缘政治推升能源价格使加息概率从25\%上调至35\%，市场定价隐含约40\%加息概率。
  \item 中国数据中心上半年智能算力新增节奏放缓，但投行维持2025-2028年20\%复合增长率，分歧在于这是供给瓶颈还是需求走弱。
\end{itemize}
\section{全球投行叙事汇编}
今日纳入 32 篇投行/券商研究，按 4 个市场、资产与产业主题系统整理。
\newpage
\subsection{AI驱动存储与光互连结构性扩张，半导体需求复苏范围扩大}
\textbf{AI基础设施投资正从GPU算力向存储、光互连和功率半导体等配套领域扩散，形成多维度、可量化的需求增长；同时汽车与工业半导体在库存出清后进入周期性复苏，但不同细分市场的增速和可持续性存在显著分化。}
\subsubsection*{主流共识}
\begin{itemize}
  \item 代理型AI将推动全球存储用量在2026-2030年间增长超过七倍，供需缺口持续存在。
  \item AI数据中心和光互连是半导体公司最清晰、最可量化的增长引擎，订单可见性显著改善。
  \item 汽车半导体需求正经历库存过低与终端需求共同推动的加速复苏，中国电动车是主要驱动力。
\end{itemize}
\subsubsection*{分歧与条件差异}
\begin{itemize}
  \item 内存接口芯片市场中，MRDIMM渗透率假设（25\%-30\% vs 保守情景）导致2030年TAM预期差异巨大（200亿美元 vs 80亿美元）。
  \item 存储供给瓶颈的持续性：部分观点认为产能扩张（30-40\% CAGR）可逐步弥合缺口，另一观点认为洁净室、设备、材料等多重瓶颈将使供需偏紧持续数年。
  \item 汽车半导体复苏的可持续性：NXP等受益于软件定义汽车架构转型（单车价值量提升），而TI/STMicro等更多依赖补库周期，后者持续性存疑。
  \item 中国DRAM制造商长鑫存储的市占率提升路径：乐观预期2028年达18\%，但美国MATCH法案等外部风险可能限制其先进节点获取。
\end{itemize}
\subsubsection*{机构观点}
\paragraph{Bernstein} 全球内存接口芯片市场2030年规模上调至约200亿美元（2025-2030年CAGR约65\%），MRDIMM升级使单模组接口硅含量达传统RDIMM的10倍，MRCD和MDB芯片将贡献73\%的市场规模。即使保守情景下（MRDIMM渗透率低于预期），市场仍可达约80亿美元。三家头部供应商（澜起、瑞萨、Rambus）合计份额超90\%，但Rambus受历史诉讼影响份额增长受限，且研发效率低于澜起。
{\small\color{refgray}数据：全球内存接口芯片TAM 2030年约200亿美元，CAGR 65\%\par}
{\small\color{refgray}数据：MRDIMM模组接口硅含量为传统RDIMM的10倍\par}
{\small\color{refgray}数据：MRCD+MDB芯片贡献2030年市场73\%\par}
{\small\color{refgray}数据：三家头部供应商2024年合计份额>90\%\par}
{\small\color{refgray}边际变化：Bernstein上调TAM预期，核心假设是MRDIMM渗透率2030年达25-30\%，此前市场普遍预期更低。\par}
\paragraph{Bernstein} 英伟达与博通通过长期协议锁定未来数年存储供应：SK海力士与英伟达签署超5000亿美元合作意向（含HBM4），三星与博通签署2000亿美元MOU（至2030年，含HBM和2nm代工）。这些协议规模占全球存储年收入（2027-2028年约1.3万亿美元）可观份额，有助于缓解市场对AI芯片供应能力的担忧。存储器件在AI基础设施中的战略地位正被重新定义，重要性甚至超过逻辑半导体。
{\small\color{refgray}数据：SK海力士与英伟达合作规模超5000亿美元\par}
{\small\color{refgray}数据：三星与博通MOU价值2000亿美元（至2030年）\par}
{\small\color{refgray}数据：SK集团另计划2500亿美元存储合作\par}
{\small\color{refgray}数据：全球存储年收入2027/2028年约1.3万亿美元\par}
{\small\color{refgray}边际变化：从短期订单转向长期框架协议，存储供应锁定成为AI基础设施投资的新常态。\par}
\paragraph{Bernstein} 英伟达以10亿美元认购Naver约5\%股份，标志着AI基础设施融资模式从单一资产负债表转向伙伴驱动的'新云'结构。英伟达作为战略股东，直接降低了市场对Naver AI工厂在需求可见性、硬件采购和长期生态支持方面的不确定性。此前Naver、英伟达和Brookfield已宣布共同研究千兆瓦级多租户AI云基础设施，英伟达入股使合作从概念走向执行。
{\small\color{refgray}数据：英伟达以1.5万亿韩元（约10亿美元）认购Naver约5\%股份\par}
{\small\color{refgray}数据：发行价每股204,500韩元，约720万股\par}
{\small\color{refgray}数据：Naver、英伟达、Brookfield计划共建千兆瓦级AI云基础设施\par}
{\small\color{refgray}边际变化：英伟达从芯片供应商转变为平台股东，利益绑定比单纯客户关系更紧密，AI基础设施融资模式出现结构性变化。\par}
\paragraph{GS} 意法半导体AI数据中心业务成为最清晰增长引擎：2026年收入预期上调至超10亿美元，2027年远高于20亿美元，主要受800G/1.6T光模块和硅光子量产驱动。订单出货比接近2倍（光互连超2倍），积压订单覆盖4.5-5个季度收入，2026年数据中心收入已被订单完全覆盖。毛利率改善路径清晰，但200mm→300mm硅生产迁移带来约60bp影响至2027年底。
{\small\color{refgray}数据：2026年数据中心收入超10亿美元，2027年远高于20亿美元\par}
{\small\color{refgray}数据：整体订单出货比接近2倍，光互连超2倍\par}
{\small\color{refgray}数据：积压订单相当于4.5-5个季度收入\par}
{\small\color{refgray}数据：制程迁移带来约60bp毛利率影响至2027年底\par}
{\small\color{refgray}边际变化：从'改善'到'可量化'，订单覆盖数据使AI需求可见性大幅提升。\par}
\paragraph{GS} 诺基亚光学业务订单急剧增长：Q2达28亿欧元（Q1为10亿欧元），AI和云客户成为核心驱动力，管理层预计约半数将在12个月内确认为收入。磷化铟晶圆和存储器等前沿光学组件供应持续紧张，客户通过签订更长期订单锁定产能。诺基亚正通过收购NXP亚利桑那工厂、扩建圣何塞和宾夕法尼亚工厂提升光学制造与封装能力（封装能力提升约十倍）。
{\small\color{refgray}数据：Q2光学订单28亿欧元，Q1为10亿欧元\par}
{\small\color{refgray}数据：约半数订单将在12个月内转化为收入\par}
{\small\color{refgray}数据：宾夕法尼亚工厂封装能力提升约十倍\par}
{\small\color{refgray}数据：2026年预计产生约8亿欧元重组费用\par}
{\small\color{refgray}边际变化：订单从10亿欧元跃升至28亿欧元，且确认节奏快于市场预期，但部分可能包含预防性采购。\par}
\paragraph{MS} 汽车半导体进入库存过低与终端需求双驱动的复苏：TI确认汽车客户库存降至'不可持续的低水平'，STMicro汽车业务Q2同比增长约13-14\%，SiC收入环比增长超30\%且book-to-bill远高于1。NXP受益于软件定义汽车架构转型（S32N、成像雷达等），加速增长驱动已占汽车收入45\%以上且保持双位数增长。数据中心对功率半导体的拉动正从概念走向可量化收入贡献。
{\small\color{refgray}数据：STMicro汽车业务Q2同比增长约13-14\%\par}
{\small\color{refgray}数据：STMicro SiC收入环比增长超30\%\par}
{\small\color{refgray}数据：NXP加速增长驱动占汽车收入45\%以上\par}
{\small\color{refgray}数据：TI汽车客户库存处于'不可持续的低水平'\par}
{\small\color{refgray}边际变化：汽车半导体复苏从一次性补库转向有持续性的周期启动，中国电动车是主要驱动力。\par}
\paragraph{NOM} 全球存储供给在未来数年难以放松，长鑫存储处于供需结构性错配窗口期。代理型AI驱动全球存储用量2026-2030年增长超七倍（即使考虑四倍内存效率折扣），而供给端受洁净室、设备、材料等多重瓶颈限制，存储制造商30-40\%的比特CAGR仍可能无法满足需求。长鑫存储比特扩张CAGR达40-45\%，全球DRAM市占率有望从约10\%提升至2028年底约18\%。
{\small\color{refgray}数据：全球存储用量2026-2030年增长超七倍\par}
{\small\color{refgray}数据：长鑫存储比特CAGR 40-45\%\par}
{\small\color{refgray}数据：全球DRAM市占率从约10\%升至2028年底约18\%\par}
{\small\color{refgray}数据：2026-2028年营收CAGR 63\%，净利润CAGR 74\%\par}
{\small\color{refgray}边际变化：NOM首次覆盖长鑫存储，强调代理型AI带来的指数级存储需求与供给瓶颈之间的结构性缺口。\par}
\paragraph{NOM} 中国大模型竞争从参数规模转向效率与系统集成：约3万亿参数是下一个里程碑，但纯参数扩展边际回报递减，竞争力取决于架构效率、训练后优化、内存管理和长周期任务表现。模型与智能体框架深度整合创造专有数据飞轮，竞争壁垒从独立模型转向模型-智能体集成系统。DeepSeek推理成本约0.8元/百万token（60\%缓存命中率），售价约1.3元，推理毛利率约40\%，新版本降价75\%后仍有利润空间。
{\small\color{refgray}数据：约3万亿参数是中国大模型下一个里程碑\par}
{\small\color{refgray}数据：DeepSeek推理成本约0.8元/百万token（60\%缓存命中率）\par}
{\small\color{refgray}数据：售价约1.3元/百万token，推理毛利率约40\%\par}
{\small\color{refgray}数据：新版本降价75\%后仍有利润空间\par}
{\small\color{refgray}边际变化：从参数规模竞赛转向效率与系统集成竞争，工程效率成为新的区分维度。\par}
\paragraph{NOM} Win Semi光学代工业务从样品阶段转向量产爬坡：光电探测器（PD）在Q2末为单一战略客户进入量产，产能已接近满载。激光产品线中VCSEL和长距离EML已量产，中距离EML预计2026下半年加入，CW激光器预计2027下半年贡献收入。光学通信收入2026年预计占高个位数百分比，2027年有望迈向双位数。毛利率改善来自产品组合优化（基础设施收入占比已与手机相当）和折旧下降至每季约7亿新台币底部。
{\small\color{refgray}数据：Q2毛利率28.2\%，营业利润率14.1\%\par}
{\small\color{refgray}数据：产能利用率从60\%提升至65\%\par}
{\small\color{refgray}数据：折旧费用每季约7.6亿新台币，接近底部\par}
{\small\color{refgray}数据：Q3毛利率指引升至低30\%区间\par}
{\small\color{refgray}边际变化：光学代工从样品阶段进入量产爬坡，单一客户产能已满载，产品线向中距离EML和CW激光器扩展。\par}
\subsubsection*{关键数据}
\begin{itemize}
  \item 全球内存接口芯片TAM 2030年约200亿美元，CAGR 65\%
  \item SK海力士与英伟达合作规模超5000亿美元
  \item 全球存储年收入2027/2028年约1.3万亿美元
  \item 意法半导体2026年数据中心收入超10亿美元，2027年远高于20亿美元
  \item 诺基亚Q2光学订单28亿欧元（Q1为10亿欧元）
  \item 长鑫存储全球DRAM市占率有望从10\%升至2028年底约18\%
  \item DeepSeek推理成本约0.8元/百万token，推理毛利率约40\%
  \item Win Semi光学代工PD已量产，产能接近满载
\end{itemize}
\begin{figure}[htbp]
\centering
\includegraphics[width=0.88\linewidth]{figures/fig_007.jpg}
\caption{EXHIBIT 2 - EXHIBIT 3: We raise our global memory interface TAM projection to USD 20Bn in 2030, with MRDIMM's MRCD and MDB contributing 73\%, driven by significantly higher chipset value per module Global memory interface chips TAM (Bn USD)}
\end{figure}
\begin{figure}[htbp]
\centering
\includegraphics[width=0.88\linewidth]{figures/fig_011.jpg}
\caption{EXHIBIT 1 - EXHIBIT 1: Samsung 2.5D Cube-S is similar to TSMC CoWoS-S, and 2.3D Cube-E \& -R are similar to TSMC CoWoS-L \& -R, respectively. EXHIBIT 2: TSMC offers three variants of CoWoS. EXHIBIT 3: Consensus memory industry to generate c.}
\end{figure}
\begin{figure}[htbp]
\centering
\includegraphics[width=0.88\linewidth]{figures/fig_006.jpg}
\caption{EXHIBIT 1 - EXHIBIT 1: Rambus has two primary business segments but reports financial statements in three lines. Two business segments ride on the wave of CPU Renaissance and AI ASIC respectively The revenue mix and growth rate numbers are o}
\end{figure}
\begin{figure}[htbp]
\centering
\includegraphics[width=0.88\linewidth]{figures/fig_008.jpg}
\caption{EXHIBIT 3 - EXHIBIT 4: The expansion of the TAM is driven by three core factors, each of which is expected to accelerate in the second half of this decade relative to the past 5 years Driver 1: Global server CPU shipment vol (Mn units)}
\end{figure}
{\small\color{refgray}本节综合 9 篇研究；涉及 Bernstein、GS、MS、NOM。}
\newpage
\subsection{消费与互联网：监管落地与结构性分化}
\textbf{中国互联网平台监管进入终局阶段，携程反垄断处罚落地消除不确定性，但竞争格局重塑开启；全球奢侈品消费受中国游客流向与区域价差驱动，品牌间分化加剧；新兴市场消费龙头通过增长业务转型对冲核心业务压力。}
\subsubsection*{主流共识}
\begin{itemize}
  \item 携程反垄断罚款金额虽高于阿里和美团案例，但公司流动性充裕（1040亿元），一次性财务影响可控，市场更关注业务调整的长期竞争格局变化。
  \item 全球奢侈品欧洲市场表现分化由三个结构性因素决定：对中国游客依赖度、区域价差大小、美国游客消费替代能力。
  \item 新兴市场消费企业（如TATACONS）通过增长业务（即饮、调味品、有机食品）实现高增速，对冲核心茶盐业务成本压力，利润率有望扩张。
\end{itemize}
\subsubsection*{分歧与条件差异}
\begin{itemize}
  \item GS认为携程处罚落地后市场反应偏正面，参考阿里和美团案例；NOM则强调排他协议取消后，美团和抖音可能缩小酒店供给差距，竞争加剧将长期压制利润率。
  \item Bernstein指出硬奢（历峰）因低中国依赖度和统一定价表现稳健，而软奢（Moncler）高依赖度叠加高价差导致欧洲收入下滑8\%；但Hermès虽高依赖度却因美国游客缓冲预期表现不同。
  \item 旺旺传统批发渠道双位数下滑与运营费用高个位数增长，估值折价（9x PE vs 同行11-14x）反映市场对其转型周期长度的分歧。
\end{itemize}
\subsubsection*{机构观点}
\paragraph{GS} 携程反垄断罚款51.79亿元（基于2025年国内收入7.5\%），高于阿里4\%和美团3\%，但公司持有1040亿元流动性，足以覆盖。监管未对抽佣率、市场份额或与同程合作施加额外限制，且携程已于1月停止违规操作。参考阿里和美团处罚公告后股价分别上涨8\%和9\%，市场可能正面解读终局信号。
{\small\color{refgray}数据：罚款51.79亿元，占2025年国内收入7.5\%\par}
{\small\color{refgray}数据：携程持有流动性1040亿元\par}
{\small\color{refgray}数据：阿里/美团罚款比例分别为4\%/3\%\par}
{\small\color{refgray}数据：阿里/美团处罚后股价分别上涨8\%/9\%\par}
{\small\color{refgray}边际变化：监管从调查期进入终局阶段，不确定性消除，且未施加额外结构性限制，市场焦点从罚款规模转向竞争格局。\par}
\paragraph{GS} 旺旺中国FY1Q26收入同比-6\%，净利润同比-38\%，低于市场预期（-5\%/-13\%）。传统批发渠道收入双位数下滑，运营费用高个位数增长。当前9x FY26E PE，低于卫龙（11x）、洽洽（13x）、盐津铺子（14x），反映市场对其转型周期长度的怀疑。
{\small\color{refgray}数据：FY1Q26收入同比-6\%，净利润同比-38\%\par}
{\small\color{refgray}数据：传统批发渠道收入双位数下滑\par}
{\small\color{refgray}数据：运营费用高个位数增长\par}
{\small\color{refgray}数据：旺旺9x PE vs 卫龙11x/洽洽13x/盐津铺子14x\par}
{\small\color{refgray}边际变化：旧渠道失效与新投入尚未见效的过渡期，估值折价隐含市场对转型周期延长的预期。\par}
\paragraph{NOM} 携程被认定在线酒店预订市场占比近60\%，罚款约53亿元（占净现金11\%），超出此前30亿元预估。排他协议取消后，美团和抖音可能缩小酒店供给差距，参考2021年案例后其他平台份额从17\%涨至40\%。预计FY26/FY27非GAAP经营利润率分别下降3pp/1pp，低于市场共识2\%和9\%。
{\small\color{refgray}数据：携程在线酒店预订市场份额近60\%\par}
{\small\color{refgray}数据：罚款53亿元，占净现金11\%\par}
{\small\color{refgray}数据：NOM此前预估30亿元\par}
{\small\color{refgray}数据：2021年案例后其他平台份额从17\%升至40\%\par}
{\small\color{refgray}边际变化：罚款金额超预期，但市场更关注排他协议取消后竞争格局的结构性变化，利润率承压预期低于共识。\par}
\paragraph{NOM} TATACONS FY1Q26印度收入同比+12\%，营业利润率13.5\%，营业利润+19.3\%。增长业务（即饮、Sampann、Capital Foods、Organic India）占印度业务36\%，同比增速47\%，管理层预期30\%+增速可持续。核心茶盐业务提价1-2\%和6\%以对冲成本上涨，全年OPM有望扩张50-75bp。
{\small\color{refgray}数据：印度收入同比+12\%，营业利润率13.5\%\par}
{\small\color{refgray}数据：增长业务占比36\%，同比增速47\%\par}
{\small\color{refgray}数据：即饮销量+35\%，Sampann+58\%，Capital Foods+40\%\par}
{\small\color{refgray}数据：茶盐提价1-2\%和6\%\par}
{\small\color{refgray}边际变化：增长业务从边缘贡献者转变为主要驱动力，结构性转型正在重塑公司DNA，利润率扩张预期强化。\par}
\paragraph{Bernstein} 全球奢侈品欧洲市场表现分化由三个因素决定：对中国游客依赖度、区域价差、美国游客替代能力。硬奢（历峰）中国客群占比20-25\%，价差小，欧洲Q2有机增长7\%；软奢（Moncler）中国客群占比35-40\%，中法价差30\%，欧洲EMEA收入同比-8\%。Hermès虽中国客群占比40\%、价差超30\%，但美国游客预期提供缓冲。
{\small\color{refgray}数据：历峰中国客群占比20-25\%，欧洲Q2有机增长7\%\par}
{\small\color{refgray}数据：Moncler中国客群占比35-40\%，中法价差30\%\par}
{\small\color{refgray}数据：Moncler欧洲EMEA收入同比-8\%\par}
{\small\color{refgray}数据：软奢平均价差约20\%\par}
{\small\color{refgray}边际变化：美国游客消费成为部分品牌的缓冲变量，而Moncler因高依赖度与高价差叠加季节性因素表现最弱。\par}
\subsubsection*{关键数据}
\begin{itemize}
  \item 携程罚款51.79亿元，占2025年国内收入7.5\%，高于阿里4\%和美团3\%
  \item 携程持有流动性1040亿元，罚款占净现金11\%
  \item 旺旺FY1Q26净利润同比-38\%，传统批发渠道双位数下滑
  \item 旺旺9x FY26E PE vs 卫龙11x/洽洽13x/盐津铺子14x
  \item 历峰欧洲Q2有机增长7\%，Moncler欧洲EMEA收入同比-8\%
  \item Moncler中法价差30\%，软奢平均20\%
  \item TATACONS增长业务占比36\%，同比增速47\%
  \item TATACONS全年OPM目标扩张50-75bp
\end{itemize}
\begin{figure}[htbp]
\centering
\includegraphics[width=0.88\linewidth]{figures/fig_014.jpg}
\caption{EXHIBIT 1 - FIGURE 1: Shopping in Europe has remained relatively attractive for Chinese consumers looking to purchase Hermès and Moncler products - even after accounting for near-shore alternatives such as Japan. More recent cross-continental}
\end{figure}
\begin{figure}[htbp]
\centering
\includegraphics[width=0.88\linewidth]{figures/fig_015.jpg}
\caption{FIGURE 1 - EXHIBIT 3: Strong demand in the USA, on the other hand, presents a potential offset for luxury brands come 2Q26E, if some of this overflows into Europe as a sequential improvement in American tourist spending. Moncler, faced with r}
\end{figure}
\begin{figure}[htbp]
\centering
\includegraphics[width=0.88\linewidth]{figures/fig_034.jpg}
\caption{Exhibit 1 - Exhibit 1: TCOM earnings breakdown Exhibit 2: Alibaba and Meituan share price during and after SAMR investigations Exhibit 4: TCOM historical forward P/E ex. exceptional}
\end{figure}
\begin{figure}[htbp]
\centering
\includegraphics[width=0.88\linewidth]{figures/fig_036.jpg}
\caption{Exhibit 3 - Exhibit 4: TCOM historical forward P/E ex. exceptional \#\# Price Target Risks and Methodology – Trip.com Group Valuation methodology: We have a Buy rating on Trip.com (TCOM/9961.HK) with 12-month target prices of US\textbackslash{}\$71/HK\textbackslash{}\$560 ba}
\end{figure}
{\small\color{refgray}本节综合 5 篇研究；涉及 Bernstein、GS、NOM。}
\newpage
\subsection{医疗与工业：技术迭代与贸易格局重塑}
\textbf{医疗领域，助听器市场因Oticon Zeal新品快速放量而重塑竞争格局，Moderna和Vertex的关键数据节点将决定其估值走向；工业领域，中国资本品出口加速渗透欧洲优势领域，三星重工估值回调后进入合理区间，但中国产能扩张仍是压制因素。}
\subsubsection*{主流共识}
\begin{itemize}
  \item 助听器市场技术迭代加速，定制式产品价值占比从15\%升至22\%，推动行业均价提升。
  \item 中国在全球资本品出口中的份额持续上升，从2005年的7\%增至2025年的24\%，欧洲份额同期从54\%降至43\%。
  \item 宁德时代电池出货量增长强劲，上半年同比增长约60\%，储能业务增速（88\%）显著高于电动汽车电池（46\%）。
\end{itemize}
\subsubsection*{分歧与条件差异}
\begin{itemize}
  \item Oticon Zeal的份额爬坡速度（两个月+420bp）是否可持续，还是仅为新品短期脉冲效应。
  \item Moderna黑色素瘤三期数据若未达预期，市场是否已充分定价下行风险（期权隐含回归基础估值约\$25）。
  \item Vertex的SION CF数据门槛（汗液氯离子降低10mmol/L）与市场预期存在差异，可能影响短期股价。
  \item 中国船厂产能扩张对新造船价格的压制程度，是否已超过三星重工F-LNG订单带来的增量收益。
\end{itemize}
\subsubsection*{机构观点}
\paragraph{GS} Oticon Zeal在VA渠道上市两个月内价值份额从18.9\%升至23.1\%（+420bp），增速远超前代产品Intent（+210bp）及竞品。6月VA渠道整体销量同比+5\%，销售额+13.8\%，定制式助听器价值占比从15\%升至22\%，反映行业从‘卖更多’向‘卖更贵’转型。
{\small\color{refgray}数据：Oticon Zeal上市两个月VA价值份额+420bp至23.1\%\par}
{\small\color{refgray}数据：VA渠道6月销量同比+5\%，销售额+13.8\%\par}
{\small\color{refgray}数据：定制式助听器价值占比从15\%升至22\%\par}
{\small\color{refgray}边际变化：Oticon Zeal的份额爬坡速度（+420bp/2个月）显著快于前代Intent（+210bp/2个月），表明新品技术突破正在加速重塑竞争格局。\par}
\paragraph{GS} Moderna股价年初至今跑赢生物科技板块61\%，主要受黑色素瘤辅助治疗三期中期数据预期驱动。当前报价已部分反映成功预期，期权市场显示成功时股价可涨49\%，失败则回归基础业务估值约\$25。2026财年收入指引同比增速不超过10\%，GS预测7.2\%。
{\small\color{refgray}数据：Moderna年初至今跑赢生物科技板块61\%\par}
{\small\color{refgray}数据：期权市场隐含成功时上涨49\%，失败回归\$25估值\par}
{\small\color{refgray}数据：2026财年收入指引同比增速≤10\%，GS预测7.2\%\par}
{\small\color{refgray}边际变化：市场焦点从新冠疫苗收入稳定性转向肿瘤疫苗数据，黑色素瘤三期中期数据成为估值核心变量。\par}
\paragraph{GS} Vertex二季度收入预计31.55亿美元，略低于共识32.19亿美元。下半年关键催化剂包括povetacicept（IgA肾病，11月30日PDUFA）、inaxaplin（APOL1肾病，2026年下半年数据）及VX-828（新一代CFTR调节剂，2026年下半年数据）。SION的CF二期数据是短期压制因素。
{\small\color{refgray}数据：Vertex二季度收入预计31.55亿美元，低于共识32.19亿美元\par}
{\small\color{refgray}数据：povetacicept PDUFA日期为11月30日\par}
{\small\color{refgray}数据：SION CF二期数据预计8月公布\par}
{\small\color{refgray}边际变化：市场对SION CF数据的评估标准与Vertex自身存在差异，后者更关注管线深度和povetacicept上市潜力。\par}
\paragraph{GS} 中国资本品全球出口份额从2005年的7\%升至2025年的24\%，欧洲从54\%降至43\%。中国对欧盟出口增速从5月的7.6\%跃升至6月的18.5\%，增速最快的品类包括钨制品（+233\%）、光纤电缆（+111\%）、船用发动机（+75\%）。欧洲在商用车发动机、乳品机械等优势领域正被中国快速追赶。
{\small\color{refgray}数据：中国资本品全球出口份额从7\%升至24\%\par}
{\small\color{refgray}数据：中国对欧盟出口增速从7.6\%跃升至18.5\%\par}
{\small\color{refgray}数据：钨制品出口增速+233\%，光纤电缆+111\%\par}
{\small\color{refgray}数据：2024年欧盟新发起33项贸易防御调查，是2023年的近3倍\par}
{\small\color{refgray}边际变化：中国对欧出口增速显著加快，且正在欧洲传统优势领域（商用车发动机、乳品机械）实现最快份额增长。\par}
\paragraph{NOM} 宁德时代二季度净利润225亿元（同比+36\%，环比+9\%），符合预期。上半年电池出货量约430-440GWh（同比+60\%），储能电池收入增速88\%远超电动汽车电池的46\%。毛利率23.2\%（环比-1.7pp），受原材料成本上涨和收入结构变化影响。公司宣布回购200-400亿元（占市值1.1-2.3\%）。
{\small\color{refgray}数据：Q2净利润225亿元，同比+36\%，环比+9\%\par}
{\small\color{refgray}数据：上半年电池出货量约430-440GWh，同比+60\%\par}
{\small\color{refgray}数据：储能电池收入增速88\%，电动汽车电池46\%\par}
{\small\color{refgray}数据：毛利率23.2\%，环比-1.7pp\par}
{\small\color{refgray}边际变化：储能业务增速接近电动汽车两倍，但毛利率受原材料成本上升和产品结构变化影响环比下降。\par}
\paragraph{NOM} 三星重工二季度营业利润3250亿韩元（同比+58.7\%），低于共识13.2\%，受250亿韩元一次性退休金支出影响。新船坞贡献1500亿韩元额外收入，但运营利润率低于造船业务均值。2026年新订单预计170亿美元（同比+115.7\%），F-LNG是主要驱动力（预计80亿美元，同比+900\%）。估值已从4月高点回调33.4\%，当前P/B约3.7倍。
{\small\color{refgray}数据：Q2营业利润3250亿韩元，低于共识13.2\%\par}
{\small\color{refgray}数据：一次性退休金支出250亿韩元\par}
{\small\color{refgray}数据：2026年新订单预计170亿美元，同比+115.7\%\par}
{\small\color{refgray}数据：F-LNG新订单预计80亿美元，同比+900\%\par}
{\small\color{refgray}边际变化：估值大幅回调后进入合理区间，但中国船厂产能扩张对新造船价格的压制仍是评级上调的主要限制。\par}
\subsubsection*{关键数据}
\begin{itemize}
  \item Oticon Zeal上市两个月VA价值份额+420bp至23.1\%
  \item 中国资本品全球出口份额从7\%升至24\%
  \item 中国对欧盟出口增速从7.6\%跃升至18.5\%
  \item 宁德时代上半年电池出货量约430-440GWh，同比+60\%
  \item 三星重工Q2营业利润低于共识13.2\%
  \item Moderna期权市场隐含成功时上涨49\%，失败回归\$25
  \item Vertex povetacicept PDUFA日期为11月30日
  \item 2024年欧盟新发起33项贸易防御调查
\end{itemize}
\begin{figure}[htbp]
\centering
\includegraphics[width=0.88\linewidth]{figures/fig_017.jpg}
\caption{Exhibit 1 - Exhibit 2: Oticon saw strong y/y value growth Oticon y/y value growth in the VA channel Exhibit 3: ...and reached all-time highs in VA value share in June 2026 (2nd month since its the launch in the VA channel in May) Oticon - Va}
\end{figure}
\begin{figure}[htbp]
\centering
\includegraphics[width=0.88\linewidth]{figures/fig_024.jpg}
\caption{Exhibit 1 - Exhibit 1: LED lighting, excavators and heat pumps have shown the most growth in Chinese exports globally (ex Russia) and greatest slowdown in imports over the last decade Snapshot of competitive intensity - China Size of bubbl}
\end{figure}
\begin{figure}[htbp]
\centering
\includegraphics[width=0.88\linewidth]{figures/fig_025.jpg}
\caption{Exhibit 1 - Exhibit 3: Engines, dairy machinery and rail brakes are the areas where Europe is most dominant vs. China in global exports markets, but the gap has been decreasing Snapshot of competitive intensity - Global, Europe/China market sh}
\end{figure}
\begin{figure}[htbp]
\centering
\includegraphics[width=0.88\linewidth]{figures/fig_027.jpg}
\caption{Exhibit 4 - Exhibit 4: China net exports to Europe have expanded significantly and regained the rising trend post covid highs China's trade balance (net export) with Europe for 43 cap goods categories, million \textbackslash{}\$ - 12m rolling ■ China Trade}
\end{figure}
{\small\color{refgray}本节综合 6 篇研究；涉及 GS、NOM。}
\newpage
\subsection{宏观与市场主题：多重叙事下的资产定价与政策博弈}
\textbf{全球市场正同时消化中国财政紧缩与贸易结构变化、AI算力需求的结构性扩张、以及地缘政治不确定性对美联储路径的扰动，资产定价在多重叙事中寻找均衡。}
\subsubsection*{主流共识}
\begin{itemize}
  \item 中国二季度GDP增速放缓至4.3\%，主因广义财政赤字率从11.0\%大幅收窄至7.0\%，土地出让收入同比骤降42\%是核心拖累。
  \item AI算力需求长期供不应求，企业token支出上限担忧缺乏数据支撑（当前月均不足11美元），且中国大模型进展强化而非削弱算力需求。
  \item 中国贸易模式转向“多出少进”，对非美发达市场出口增长约20\%，已使当地商品价格平均下降0.7\%，日本（-1.1\%）和欧元区（-1.0\%）受影响最大。
\end{itemize}
\subsubsection*{分歧与条件差异}
\begin{itemize}
  \item 美联储7月会议路径分歧：通胀改善（6月核心CPI环比-0.02\%）支持按兵不动，但地缘政治推升能源价格使加息概率从25\%上调至35\%，市场定价隐含约40\%加息概率。
  \item 中国数据中心上半年智能算力新增节奏放缓（595 EFLOPS vs 2H25的802 EFLOPS），但GS维持2025-2028年20\%复合增长率，分歧在于这是供给瓶颈还是需求走弱。
  \item HDD行业结构性短缺vs市场担忧：MS认为每年300-400EB供应缺口持续至2028年，定价谈判已上移至25-30美元/TB，而市场对产能过剩和NAND替代的担忧被高估。
  \item 人民币中间价方向：NOM模型预测值6.7685较前次下移254个基点，主要受美元走弱驱动，但逆周期因子计入后仍偏强，政策意图与模型输出存在差异。
\end{itemize}
\subsubsection*{机构观点}
\paragraph{GS} 中国二季度GDP增速放缓至4.3\%的主因是广义财政赤字率从11.0\%收窄至7.0\%，土地出让收入同比骤降42\%是核心拖累。预计7月政治局会议将释放更强宽松信号，以推动下半年地方项目加速落地。7月制造业PMI可能跌破荣枯线至49.9。
{\small\color{refgray}数据：6月预算内财政支出增速4.0\%，远低于收入增速8.7\%\par}
{\small\color{refgray}数据：广义财政赤字率从2025年11.0\%收窄至7.0\%\par}
{\small\color{refgray}数据：二季度GDP增速4.3\%，低于全年目标区间4.5\%-5\%下沿\par}
{\small\color{refgray}数据：7月制造业PMI预测49.9，低于6月的50.3\par}
{\small\color{refgray}边际变化：此前市场将增长放缓归因于需求不足，GS明确指出财政紧缩是主因，且紧缩力度超预期。\par}
\paragraph{GS} 中国数据中心上半年智能算力新增约595 EFLOPS，低于2H25的802 EFLOPS，主因芯片供应瓶颈而非需求走弱。AI开源模型推理效率突破推动需求，GS维持2025-2028年20\%复合增长率，预计总容量从2025年16GW增至2028年27GW。华为Atlas 950 SuperPoD与英伟达B300差距从1/8缩小至1/4，但单位算力资本开支为B300的3.2倍。
{\small\color{refgray}数据：1H26新增智能算力595 EFLOPS，低于2H25的802 EFLOPS\par}
{\small\color{refgray}数据：2025-2028年数据中心需求CAGR 20\%，总容量从16GW增至27GW\par}
{\small\color{refgray}数据：华为Atlas 950 SuperPoD与B300差距从1/8缩小至1/4\par}
{\small\color{refgray}数据：国产替代单位算力资本开支为B300的3.2倍\par}
{\small\color{refgray}边际变化：此前市场担忧需求走弱，GS指出节奏放缓是供给侧阶段性约束，下半年有望重新加速。\par}
\paragraph{GS} 中国贸易模式转向“多出少进”，对非美发达市场出口增长约20\%，进口增长持续低于疫情前趋势。GS测算中国出口渗透率每提升1个百分点，接收国商品价格下降0.5\%，已使非美发达市场商品价格平均下降0.7\%，日本（-1.1\%）和欧元区（-1.0\%）受影响最大。
{\small\color{refgray}数据：中国对非美发达市场名义出口增长约20\%\par}
{\small\color{refgray}数据：中国出口渗透率每提升1pp，商品价格下降0.5\%\par}
{\small\color{refgray}数据：非美发达市场商品价格平均下降0.7\%\par}
{\small\color{refgray}数据：日本受影响最大（-1.1\%），欧元区次之（-1.0\%）\par}
{\small\color{refgray}边际变化：此前市场关注中国出口对全球通胀的推升作用，GS量化了出口增加和进口减少的双重通缩效应。\par}
\paragraph{GS} 7月FOMC会议前通胀改善（6月核心CPI环比-0.02\%），但地缘政治不确定性推升能源价格，使利率路径更加不确定。GS维持年底前按兵不动基准预测，但将最终加息概率从25\%上调至35\%。市场定价隐含约40\%加息概率，无论结果如何都将是近几十年最大意外之一。
{\small\color{refgray}数据：6月核心CPI环比-0.02\%\par}
{\small\color{refgray}数据：GS将最终加息概率从25\%上调至35\%\par}
{\small\color{refgray}数据：市场定价隐含约40\%加息概率\par}
{\small\color{refgray}数据：BEA方法论调整有望使8月同比通胀率下降约0.2个百分点\par}
{\small\color{refgray}边际变化：此前市场聚焦通胀改善，GS指出地缘政治对加息辩论的影响超过油价传导数学计算，增加了供给冲击反复的不确定性。\par}
\paragraph{GS} 截至7月23日当周，中国至美国重载集装箱船数量同比下降14.5\%，TEU货运量同比下降23\%，连续四周环比下降。中国大陆降幅（-23\%）大于亚洲其他地区（-9\%），反映供应链向东南亚转移趋势。洛杉矶港下周到港量预计环比增加2\%，但后续一周可能转为同比下降10\%。
{\small\color{refgray}数据：中国至美国重载集装箱船数量同比-14.5\%\par}
{\small\color{refgray}数据：TEU货运量同比-23\%，连续四周环比下降\par}
{\small\color{refgray}数据：中国大陆TEU降幅-23\%，亚洲其他地区-9\%\par}
{\small\color{refgray}数据：洛杉矶港下周到港量预计环比+2\%，后续一周同比-10\%\par}
{\small\color{refgray}边际变化：此前市场关注关税对贸易流量的一次性冲击，GS提供高频追踪框架显示供应链转移和补库节奏正在演变。\par}
\paragraph{GS} 沃达丰2027财年第一季度温和超预期，集团将全年EBITDA和自由现金流指引上调至区间上限，对应同比分别增长约7\%和20\%。主要驱动力来自非洲业务（Vodacom有机服务收入增长12.6\%），欧洲有机服务收入增长仍停滞在0.9\%，德国增长靠宽带ARPU和B2B数字服务支撑，但用户持续净流失。
{\small\color{refgray}数据：集团EBITDA和自由现金流指引上调至区间上限，对应同比+7\%和+20\%\par}
{\small\color{refgray}数据：Vodacom有机服务收入增长12.6\%，高于预期的11.5\%\par}
{\small\color{refgray}数据：欧洲有机服务收入增长0.9\%，与上季度0.8\%基本持平\par}
{\small\color{refgray}数据：德国宽带净增用户-9.8万，后付费用户净增-8.5万\par}
{\small\color{refgray}边际变化：此前市场关注欧洲业务改善，GS指出增量主要来自非洲周期性因素，欧洲基本面增长动能未变。\par}
\paragraph{MS} AI基础设施板块近期回调是技术性调整而非基本面转向，算力需求将在未来多年持续显著超过供给。企业token支出上限担忧缺乏数据支撑（月均不足11美元），AI使用带来的劳动力成本节约约55美元，token成本仅2-3美元。中国大模型进展强化“杰文斯悖论”，更高效模型降低使用成本反而推动更多算力需求。
{\small\color{refgray}数据：企业用户月均token花费不足11美元\par}
{\small\color{refgray}数据：AI替代人工可节省约55美元/任务，token成本仅2-3美元\par}
{\small\color{refgray}数据：前沿模型已在超过50\%人类知识任务中达到或超越人类水平\par}
{\small\color{refgray}数据：DeepMind AlphaProof框架以数百美元推理成本解决2道56年未解数学问题\par}
{\small\color{refgray}边际变化：此前市场担忧AI基建回调反映需求见顶，MS指出技术性因素是主因，算力供不应求格局将延续多年。\par}
\paragraph{MS} HDD行业处于结构性供应短缺，每年300-400EB供应缺口持续至2028年，需求可见度已延伸至2032年。定价谈判上移至25-30美元/TB，现货溢价30-40\%以上。市场对产能扩张和NAND替代的担忧被高估，HDD相对eSSD仍保持17-20倍成本优势。
{\small\color{refgray}数据：每年300-400EB供应缺口持续至2028年\par}
{\small\color{refgray}数据：2027-2028年定价谈判区间25-30美元/TB\par}
{\small\color{refgray}数据：现货EB交易溢价30-40\%以上\par}
{\small\color{refgray}数据：HDD相对eSSD保持17-20倍成本优势\par}
{\small\color{refgray}边际变化：此前市场担忧HDD定价持续性和盈利能力，MS指出产业实际信号与市场担忧存在明显偏离，定价是盈利超预期的主要驱动力。\par}
\paragraph{NOM} NOM不含逆周期因子的USD/CNY中间价模型预测值为6.7685，较前次下移254个基点，较当日官方收盘价低59个基点。预测下移主要受美元走弱驱动，美元指数贡献最大负向权重。计入逆周期因子后预测值为6.7741，较前次中间价低198个基点，幅度在收窄。
{\small\color{refgray}数据：不含逆周期因子模型预测值6.7685，较前次下移254个基点\par}
{\small\color{refgray}数据：较当日官方收盘价低59个基点\par}
{\small\color{refgray}数据：计入逆周期因子后预测值6.7741，较前次中间价低198个基点\par}
{\small\color{refgray}数据：美元指数贡献最大负向权重\par}
{\small\color{refgray}边际变化：此前市场关注中国央行主动引导人民币升值，NOM模型显示预测下移主要跟随全球美元定价，政策意图与模型输出需区分。\par}
\paragraph{Bernstein} 全球储能产业从单一成本竞争转向技术专利、供应链合规和区域化布局的多维博弈。北美市场焦点转向知识产权保护和供应链“去中国化”，LG Energy Solution起诉中国电池制造商EVE Energy涉及五项专利。中国目标到2030年可再生能源消费量提升53\%，储能支撑的可再生能源需满足20\%峰值电力需求。韩国EcoPro BM硫化物固态电解质已通过认证，预计2027年量产。
{\small\color{refgray}数据：LG Energy Solution起诉EVE Energy涉及5项专利\par}
{\small\color{refgray}数据：中国目标2030年可再生能源消费量提升53\%\par}
{\small\color{refgray}数据：储能支撑的可再生能源需满足20\%峰值电力需求\par}
{\small\color{refgray}数据：EcoPro BM固态电解质预计2027年量产\par}
{\small\color{refgray}边际变化：此前市场聚焦储能成本竞争，Bernstein指出专利诉讼和供应链合规正成为新的竞争壁垒。\par}
\paragraph{Bernstein} 量化模型与基本面研究结合选股效果显著优于单一方法。自2013年以来，同时进入BernsteinAlpha模型最优两个五分位且被分析师给予“跑赢”评级的公司，年化超额回报达9.4\%；收紧至最优五分位后升至13.2\%。自2020年起发布的12期实盘组合在等权重口径下均跑赢MSCI亚洲指数。
{\small\color{refgray}数据：量化+基本面组合年化超额回报9.4\%\par}
{\small\color{refgray}数据：最优五分位年化超额回报13.2\%\par}
{\small\color{refgray}数据：过去20年最优五分位年化超额回报6.1\%，信息比率1.35\par}
{\small\color{refgray}数据：12期实盘组合等权重口径均跑赢MSCI亚洲指数\par}
{\small\color{refgray}边际变化：此前市场对量化策略有效性存在争议，Bernstein提供连续12期实盘跟踪记录验证策略长期胜率。\par}
\paragraph{GS} 2024年7月第三至第四周，GS分析师在亚洲（香港、新加坡）与约55名读者进行了为期八天的交流。此时正值市场对AI需求和竞争格局前景担忧加剧、市场波动性较高的阶段。但报告显示，多数读者对半导体及AI相关需求和基本面仍持正面态度，尤其对电子材料领域的兴趣保持高位。读者关注的核心问题，是相关公司能否在AI材料供需趋紧的环境下及时进行资本开支并行使定价权。讨论热度...
{\small\color{refgray}数据：读者对日东纺的悲观情绪集中体现在T-glass竞争预期上\par}
{\small\color{refgray}数据：报告指出，过去三个月左右，许多读者因日东纺主力产品T-glass供需平衡出现变化以及相对较高的估值，选择做空或低配该股。市场普遍认为，公司在低热膨胀玻璃和低介电玻璃领域面临来自台湾、日本及中国大陆新进入者的竞争，其竞争优势可能被侵蚀。然而，GS近期将日东纺评级从中性上调至观点偏积极，并认为尽管有新进入者，但竞争对手产品的实际采用进度显著延迟，公司竞争优势短期...\par}
\subsubsection*{关键数据}
\begin{itemize}
  \item 中国广义财政赤字率从11.0\%收窄至7.0\%，土地出让收入同比-42\%
  \item 中国对非美发达市场出口+20\%，已使当地商品价格平均-0.7\%
  \item 6月核心CPI环比-0.02\%，GS将最终加息概率从25\%上调至35\%
  \item 中国1H26新增智能算力595 EFLOPS，低于2H25的802 EFLOPS
  \item 企业用户月均token花费不足11美元，AI替代人工可节省55美元/任务
  \item HDD每年300-400EB供应缺口持续至2028年，定价谈判上移至25-30美元/TB
  \item NOM模型预测USD/CNY中间价6.7685，较前次下移254个基点
  \item 量化+基本面组合年化超额回报13.2\%，12期实盘组合均跑赢基准
\end{itemize}
\begin{figure}[htbp]
\centering
\includegraphics[width=0.88\linewidth]{figures/fig_019.jpg}
\caption{Exhibit 1 - Exhibit 3: China's \textbackslash{}\textasciitilde{}600 EFLOPS intelligent computational power addition in 1H26 tracks below 2H25 run rate, yet we expect the pace to re-accelerate in 2H26E}
\end{figure}
\begin{figure}[htbp]
\centering
\includegraphics[width=0.88\linewidth]{figures/fig_029.jpg}
\caption{Exhibit 1 - Exhibit 2: Chinese Import Growth Has Undershot Its Pre-Pandemic Trend}
\end{figure}
\begin{figure}[htbp]
\centering
\includegraphics[width=0.88\linewidth]{figures/fig_037.jpg}
\caption{Exhibit 1 - Exhibit 1: Re-escalation of the War with Iran Has Pushed Oil and Refined Product Prices Higher Again We expect the FOMC to leave the fed funds rate unchanged at its July meeting next week. President Logan has expressed support fo}
\end{figure}
\begin{figure}[htbp]
\centering
\includegraphics[width=0.88\linewidth]{figures/fig_047.jpg}
\caption{Exhibit 2 - Exhibit 2: On average, these names have been underperformers over the past month, despite strong fundamentals and AI infrastructure exposure. AI Infrastructure "Ways to Play" List: Average Stock Performance (Past 1 Month) \#\# What}
\end{figure}
{\small\color{refgray}本节综合 12 篇研究；涉及 Bernstein、GS、MS、NOM。}
\newpage
\section{辅助信号｜战略咨询}
本板块聚焦波士顿咨询两项研究：AI优先组织设计如何从试点走向规模化价值，以及韩国股市重估的结构性分化与政策驱动逻辑。
\subsection{为AI重新设计公司：从试点到规模化价值}
\textbf{多数企业AI试点效果有限，根源在于组织模式仍为纯人类劳动力设计；真正实现规模化价值的5\%企业，通过围绕智能体系统重构工作结构、治理与流程，将AI从工具升级为运营核心。}
\begin{itemize}
  \item BCG研究显示，仅5\%的企业从AI中获得可规模化价值，核心障碍不是技术，而是组织模式未适配AI——传统流程围绕人类角色和层级设计，导致试点碎片化、所有权不清、规模化受限。
  \item 欧洲能源企业案例：放弃将AI插入现有工作流，转而按AI优先模式重新设计端到端客户旅程，结果满意度迅速匹配并有望超越人类表现，对外部服务提供商依赖度降低90\%以上，释放现金流用于加速转型。
  \item AI优先组织的工作围绕智能体系统动态协调，人类负责定义目标、设定约束、管理不确定性，AI在参数范围内自主寻优；这种模式要求同步调整流程、数据、系统及治理结构。
  \item 该企业采用“自筹资金”策略，优先改造能快速创造价值的客户旅程，前三个月即产生运营成本节省，验证了“先验证、再扩展”的转型节奏。
\end{itemize}
\subsection{韩国股市重估：结构性分化与政策驱动}
\textbf{2025年KOSPI TSR约76\%，但上涨高度集中于半导体、造船、国防和核电四个行业；超过60\%上市公司仍以低于账面价值交易，下一阶段重估取决于企业能否将战略重心转向价值最大化。}
\begin{itemize}
  \item 2025年KOSPI总股东回报率约76\%，远超全球前十大指数均值22\%，其中75个百分点来自资本利得，股息仅贡献1个百分点；PBR从2024年0.8倍升至1.4倍，2026年预期进一步升至1.9倍。
  \item 剔除半导体、造船、国防和核电四个关键行业后，2026年预期PBR降至1.0倍；半导体行业ROE预期达60\%，市值占比从29\%跃升至52\%，而其余行业平均ROE仅约7\%，仍低于约10\%的股权成本。
  \item 第一轮重估由两股力量驱动：四个关键行业的盈利复苏（AI存储芯片需求、全球造船/国防订单周期、核电重启）以及新政府资本市场振兴政策（商法修订、价值提升计划、股息分离课税、禁止新股低价发行）。
  \item 政策层面系统性推进正在改变市场对韩国结构性折价的预期，但价值提升披露参与率仅约28\%，远低于日本经验；下一阶段重估取决于企业能否将战略重心转向企业价值最大化，并同步调整资本配置、股东回报和市场沟通。
\end{itemize}
\begin{figure}[htbp]
\centering
\includegraphics[width=0.88\linewidth]{figures/fig_069.jpg}
\caption{Exhibit 1 - KOSPI: \#1 Index Return Globally In 2025, the Korean stock market delivered an overwhelmingly strong performance relative to major global markets. The KOSPI's TSR was approximately 76\%, significantly exceeding the average of the top 10 global indices (approxima}
\end{figure}
\begin{figure}[htbp]
\centering
\includegraphics[width=0.88\linewidth]{figures/fig_072.jpg}
\caption{Exhibit 3 - \#\# 1. Market Re-rating Driven by Earnings Growth in Four Key Sectors: Semiconductors, Shipbuilding, Defense, and Nuclear Power The earnings improvement in the four key sectors drove the market's re-rating. The total KOSPI market cap expanded approximately thre}
\end{figure}
\newpage
\section{辅助信号｜智库/国际机构}
经合组织两份报告分别聚焦医疗数据互操作性的经济价值与障碍，以及G7国家AI劳动力市场政策的现状与缺口。前者强调治理与信任是比技术更关键的瓶颈，后者指出技能培训偏重技术操作而忽视AI素养，且多数国家缺乏针对AI的专门法规。
\subsection{医疗数据互操作性的价值与结构性障碍}
\textbf{医疗数据互操作性可释放相当于卫生支出2.7\%至6.6\%的宏观价值，但技术标准进展未能解决治理碎片化、信任缺失和运营瓶颈，后者才是实现价值的关键短板。}
\begin{itemize}
  \item 经合组织估算，加拿大、芬兰和英国等国的实证显示，有效的数据互操作性每年可释放相当于卫生支出2.7\%至6.6\%的宏观环境价值，来源包括服务交付优化、错误减少和更及时的治疗。
  \item 当前结构性障碍显著：81.5\%的医生认为互操作性不足对患者安全构成潜在不确定性，患者因重复提供信息而对医疗系统的信任度下降15\%，诊断错误占医疗成本的约17.5\%，其中相当比例与数据互操作不畅相关。
  \item 技术标准已有进展——所有受访国家都在使用健康数据标准，41\%完全依赖国际标准，识别出28种现行标准——但标准采用不等于互操作性实现。
  \item 受访者提及最多的障碍包括：不一致的数据共享实践和去标识化做法（70.59\%）、区域间基础设施和技术碎片化（66.67\%），以及碎片化的治理和机构协调（60.78\%）。
  \item 报告核心判断：技术层面的进展并未解决法律和运营层面的根本瓶颈，真正的挑战在于人、治理和领导力。
\end{itemize}
\begin{figure}[htbp]
\centering
\includegraphics[width=0.88\linewidth]{figures/fig_062.jpg}
\caption{Figure 1 - \#\# What is interoperability in healthcare? 18. Interoperability is the ability of different health systems, devices, applications, or organisations to securely exchange, interpret, and use data in a coordinated and meaningful way, with minimal human interventi}
\end{figure}
\begin{figure}[htbp]
\centering
\includegraphics[width=0.88\linewidth]{figures/fig_065.jpg}
\caption{Figure 4 - 30. This analysis is guided by the Interoperability Tower (see Figure 4) which provides a frame of reference across local, regional, national, and cross-border levels of data exchange. With the complex nature of interoperability, and its common challenges invo}
\end{figure}
\subsection{AI劳动力市场政策：技能培训密集但偏技术，法规框架仍依赖旧体系}
\textbf{G7国家AI劳动力市场政策集中在技能培训，但偏重技术操作而非AI素养；多数国家未为AI单独设立政策，而是纳入现有劳动法规，在隐私、透明度和问责方面政策仍处早期阶段。}
\begin{itemize}
  \item 所有被调查国家都提供再培训和技能提升项目，但大多数聚焦于技术技能，而非通用AI素养——工人对AI如何影响自身工作、如何与AI系统协作的理解未被纳入主流培训目标。
  \item 针对因AI而失业的工人，各国主要依赖标准就业支持政策和社会保护，而非专门设计的AI转型支持；对于仍在职但面临AI冲击不确定性较高的工人，针对性措施仍然稀缺。
  \item 在隐私和数据保护方面，各国首先依赖现有框架，许多国家提供额外指南和工具帮助雇主理解义务，但数据保护和隐私框架对工人及AI系统的具体影响并不总是被清晰界定或向雇主充分解释。
  \item 报告提出评估AI劳动力市场政策成熟度的三个维度：是否有专门针对AI的法规而非仅依赖旧框架、技能培训是否涵盖AI素养而非仅技术操作、社会对话中是否建立了AI专业知识基础。目前多数国家在这三个维度上仍有明显提升空间。
  \item 社会对话和集体谈判可以帮助解决AI在劳动力市场使用中的实际挑战，但前提是各方具备相应的技术理解能力——当前谈判方在评估AI影响时可能缺乏足够技术知识。
\end{itemize}
\begin{figure}[htbp]
\centering
\includegraphics[width=0.88\linewidth]{figures/fig_067.jpg}
\caption{Figure 1 - At the same time, AI has the potential to help people with disabilities operate on the labour market. For instance, Touzet (2023[13]) interviewed over 70 stakeholders and identified 142 examples of AI-powered solutions to support people with disability, and em}
\end{figure}
\begin{figure}[htbp]
\centering
\includegraphics[width=0.88\linewidth]{figures/fig_068.jpg}
\caption{Figure 2 - \# 4 Non-discrimination in the labour market AI systems can help identify and address inequalities in the workplace by providing quantitative insights and supporting fairer decision making. For example, in the OECD AI surveys, 57\% of workers in manufacturing an}
\end{figure}
\section{结语}
后续需验证的关键节点包括：Moderna黑色素瘤三期数据是否达预期、Vertex的SION CF数据门槛与市场预期差异、中国数据中心智能算力新增节奏放缓的持续性、以及美联储7月会议路径在地缘政治扰动下的最终选择。
\newpage
\section{覆盖概览}
投行/券商 32 篇；战略咨询 2 篇；智库/国际机构 2 篇。\par
投行覆盖：GS 15篇 / NOM 8篇 / Bernstein 6篇 / MS 3篇\par
\section{Disclaimer}
{\scriptsize\color{refgray}
This community is a paid learning and information-sharing community created for general educational purposes only. The membership fee is charged solely for access to community discussions, educational content, organization, and operational costs. It is not an advisory fee, management fee, performance fee, brokerage fee, or compensation for personalized financial, investment, legal, tax, or professional advice.\par
I participate and share information solely as an individual and for educational discussion among community members. I am not acting as a financial adviser, investment adviser, broker, dealer, portfolio manager, legal adviser, tax adviser, accountant, fiduciary, or any other licensed professional adviser. Nothing shared in this community constitutes, and should not be interpreted as, financial advice, investment advice, legal advice, tax advice, a recommendation, solicitation, offer, endorsement, instruction, or guarantee to buy, sell, hold, short, trade, or invest in any security, cryptocurrency, fund, derivative, strategy, or other financial product.\par
All content shared in this community, including messages, discussions, links, files, charts, examples, market commentary, watchlists, AI-generated or AI-polished summaries, and third-party materials, is general, impersonal, and educational in nature. It does not take into account any individual's financial situation, investment objectives, risk tolerance, investment horizon, liquidity needs, tax status, legal restrictions, or suitability requirements. No content should be relied upon as suitable for any specific person, account, portfolio, or financial decision.\par
Information shared in this community may be incomplete, inaccurate, outdated, speculative, AI-generated, AI-edited, or based on third-party internet sources that have not been independently verified. I make no representation or warranty as to the accuracy, completeness, timeliness, reliability, or suitability of any information shared.\par
Financial markets involve substantial risk, including the possible loss of principal. Past performance, historical data, hypothetical examples, simulated results, backtests, personal opinions, or educational examples do not guarantee future results. Each member is solely responsible for conducting their own research, exercising independent judgment, and making their own decisions. Before making any financial, investment, legal, or tax decision, members should consult qualified and licensed professionals.\par
Participation in this community, payment of any membership fee, or communication with me or other members does not create any adviser-client, fiduciary, professional, contractual, agency, partnership, employment, or other advisory relationship. By joining or participating in this community, you acknowledge and agree that you are solely responsible for your own decisions, actions, transactions, profits, losses, risks, and consequences, and that I shall not be liable for any direct, indirect, incidental, consequential, financial, legal, tax, or other loss or damage arising from or related to any content, discussion, information, or material shared in this community.\par
}
\newpage
\begin{center}
{\Large 更多详情报告kcdesk.com}\par\vspace{0.8cm}
\includegraphics[width=0.58\linewidth]{\detokenize{../../prompts/zsxq_img.jpg}}
\end{center}
\end{document}
